\documentclass[11pt]{article}
\usepackage[a4paper,margin=24mm]{geometry}
\usepackage[T1]{fontenc}
\usepackage{lmodern,amsmath,amssymb,amsthm,mathtools,booktabs,microtype}
\usepackage[colorlinks=true,linkcolor=blue,urlcolor=blue,citecolor=blue]{hyperref}
\newtheorem{proposition}{Proposition}
\newcommand{\eps}{\epsilon}
\newcommand{\ket}[1]{|#1\rangle}
\newcommand{\B}{\widehat B}
\newcommand{\PP}{\widehat P}
\newcommand{\Bt}{B_{\rm tr}}
\newcommand{\U}{\mathcal U}
\newcommand{\F}{\mathcal F}
\newcommand{\D}{\mathcal D}
\title{An orthogonal frame for the matter chord module\\
\large Moving-core symmetry limits and a positive null candidate}
\author{Alan Day \and Chengheng\thanks{AI research assistant (OpenAI).}}
\date{8 November 2026\quad Local research draft}
\begin{document}\maketitle
\begin{abstract}
We give an orthogonal realization of the established one-matter
chord representation, separating a left q-oscillator from its
vacuum multiplicity. An ordinary displacement on the first factor
then defines an explicit moving unitary frame on the entire
module. In a joint large-displacement, small-deformation limit,
every fixed finite word of the Lin--Stanford generators and the
Hamiltonian boost converges on the finite-support core to the
positive discrete series of $\mathfrak{sl}_2$. This extends the
preceding endpoint single-vector comparison to descendants and
commutator norms. The combination $E-\widehat P$ has a positive,
boost-covariant limit with the standard Laguerre spectral measure.
Finally, an operator-norm bound on the neutral field controls a
specified insertion of the full selected density without replacing
it by an energy-only density. The limit retains a null-generator
candidate but does not determine a mass or a multiparticle law.
\end{abstract}

\section{Setting and existing constructions}
Work after the established fixed-$\lambda$, half-filled double
scaling, with $q=e^{-\eps}$, $\eps=4\lambda$,
$K=H/(2J_s)=h/\sqrt\eps$ and $\hbar=1$.
The one-insertion module has weight $d>0$, crossing $r=q^d$,
and seed $\Omega_d=\ket{d;0,0}$. The actual charged curvature
channel has $d=2$. This module is the existing tracial one, not
a new independent-probe model.

Lin--Stanford (LS) supply its ladder algebra and length-preserving
generators \cite{LS}. Okuyama's doubled-space map and Xu's
isometric formulation explicitly involve the matter-dependent
inner product \cite{Oku,Xu}. A linear tensor map in those
coordinates cannot simply be treated as a product-inner-product
unitary. We instead orthogonalize using the kernel of $a_L$.
This is an elementary realization of the same representation;
neither the factorization idea nor the resulting discrete series
is claimed as a new representation-theoretic discovery.

We retain the earlier conventions
$J_{ij}=a_i^\dagger a_j-[\bar n]$, $c=q^{\bar n/2}$,
$E=-(J_{LL}+J_{RR})/(2c)$,
$B=(J_{LL}-J_{RR})/(2c)$,
$P=i(J_{LR}-J_{RL})/2$, and
$\B=B/\sqrt{1-c^2}$, $\PP=P/\sqrt{1-c^2}$.
The rescaling acts separately on each total level.
The physical comparison is $\Bt=(K_L-K_R)/2$.

\newpage
\section{An exact orthonormal tensor frame}
Write $[n]=(1-q^n)/(1-q)$ and
$(r^2;q)_k=\prod_{j=0}^{k-1}(1-r^2q^j)$.
Define the normalized left-vacuum vectors
\begin{equation}
 \xi_k=\frac{\prod_{j=0}^{k-1}
       (a_R^\dagger-rq^j a_L^\dagger)\Omega_d}
                {\sqrt{[k]!(r^2;q)_k}},\qquad
 e_{n,k}=\frac{(a_L^\dagger)^n}{\sqrt{[n]!}}\xi_k.
 \label{frame}
\end{equation}
All factors in the numerator commute. Empty products are one.

\begin{proposition}\label{orthogonal}
The $e_{n,k}$, $n,k\ge0$, form an orthonormal basis of the
one-matter module. The unitary
$\F_{q,d}:\ell^2(\mathbb N_0)\otimes\ell^2(\mathbb N_0)
\to\mathcal H_d$, defined by $\ket{n}\otimes\ket{k}\mapsto e_{n,k}$,
gives
\begin{equation}
 a_L=a_q\otimes1,\qquad
 a_R=a_q\otimes R+1\otimes b_q,\qquad
 \bar n=N+K_0+d,\qquad R=q^{d+K_0},\label{tensor}
\end{equation}
where tensor symbols are suppressed on the right in $\bar n$,
and
\[
 a_q\ket n=\sqrt{[n]}\ket{n-1},\quad
 b_q\ket k=\sqrt{[k](1-r^2q^{k-1})}\ket{k-1},\quad K_0\ket k=k\ket k.
\]
\end{proposition}
\begin{proof}
Suppose $\xi_k$ is a normalized $a_L$-vacuum at total level $k$.
The relations $[a_L,a_R^\dagger]=q^{\bar n}$ and
$a_La_L^\dagger-qa_L^\dagger a_L=1$ show that
$(a_R^\dagger-rq^ka_L^\dagger)\xi_k$ is again a left vacuum.
The q-oscillator relation gives the ladder norm $[n]!$.
Different $n$ at the same total level have distinct eigenvalues
$[n]$ of the self-adjoint $a_L^\dagger a_L$, hence are orthogonal.
Inductively there is exactly one new left-vacuum direction per
level, because the level-$s$ space has dimension $s+1$.

Let $s_k$ be the norm of the new vacuum vector. Since $a_R$
commutes with $a_L$, it maps the vacuum line at level $k$ to
the previous one; adjointness gives $a_R\xi_k=s_{k-1}\xi_{k-1}$.
The right q-oscillator relation yields
\[
 s_k^2=1+q s_{k-1}^2-r^2q^{2k}
      =[k+1](1-r^2q^k),\qquad s_{-1}=0.
\]
This proves the normalization in \eqref{frame} and positivity
of each new direction. Commuting $a_R^\dagger$ through
$(a_L^\dagger)^n$ gives \eqref{tensor}; the annihilators follow
by adjointness. Each total level now has a complete orthonormal
set of $s+1$ vectors.
\end{proof}
For $k=0$, $e_{n,0}$ is the preceding endpoint basis; its
orthogonal $f_n$ is exactly $e_{n-1,1}$.
Diagonalizing $a_q+a_q^\dagger$ makes the right Hamiltonian
an Al-Salam--Chihara Jacobi operator with parameters
$re^{\pm i\theta_L}$, consistent with the known conditional
matter measure \cite{Xu,DLMFq}.

\newpage
\section{The moving frame and its limiting generators}
Let $a_1$ be the ordinary oscillator on the first factor, and
set $D(-t)=\exp[-t(a_1^\dagger-a_1)]$, $t=\sqrt m$.
This is an ordinary unitary displacement, not q-thermal
evolution. Define
\begin{equation}
 \U_{q,d,m}=\F_{q,d}(D(-\sqrt m)\otimes1),\qquad
 \D=\operatorname{span}\{\ket j\otimes\ket k:j,k<\infty\}.
 \label{moving}
\end{equation}
Its ground vector is the moving coherent state used to compare
with the true heat vector in the preceding note.

On the second factor introduce the standard discrete-series
operators
\begin{align}
 L_-\ket k&=\sqrt{k(2d+k-1)}\ket{k-1},&L_+&=L_-^\dagger,\\
 E_0&=d+K_0,&B_0&=(L_++L_-)/2,&P_0&=i(L_+-L_-)/2.
\end{align}
They act as the identity on the first factor. Their commutators
are $[B_0,P_0]=iE_0$, $[E_0,P_0]=iB_0$,
$[B_0,E_0]=iP_0$, with
$E_0^2-B_0^2-P_0^2=d(d-1)$.

\begin{proposition}\label{words}
Take $m\to\infty$, $\eps\to0$ with $\eps m\to0$ and fixed
$d>0$. For every fixed finite noncommutative polynomial $p$
and every $v\in\D$,
\begin{equation}
 \|p(E,\B,\PP,\Bt)\U v
 -\U p(E_0,B_0,P_0,-B_0)v\|\longrightarrow0.\label{wordlimit}
\end{equation}
In particular the three $\mathfrak{sl}_2$ defects and the
Casimir defect converge in norm on this moving core.
\end{proposition}

Here are the coefficients used to prove it. Put
\[
R_k=q^{d+k},\qquad c_{n,k}=q^{(n+k+d)/2},\qquad
s_k=\sqrt{[k+1](1-r^2q^k)},\qquad s_{-1}=0.
\]
The diagonal coefficients are
\begin{align}
 E_{n,k;n,k}&=
 \frac{[d+k]-(1-R_k)^2[n]/2-s_{k-1}^2/2}{c_{n,k}},\\
 \B_{n,k;n,k}&=
 \frac{(1-R_k^2)[n]-s_{k-1}^2}
 {2c_{n,k}\sqrt{1-c_{n,k}^2}},\qquad \PP_{n,k;n,k}=0.
\end{align}
The edge $(n,k)\to(n-1,k+1)$ has coefficients
\begin{equation}
 -\frac{R_ks_k\sqrt{[n]}}{2c_{n,k}},\qquad
 -\frac{R_ks_k\sqrt{[n]}}{2c_{n,k}\sqrt{1-c_{n,k}^2}},\qquad
 -\frac{i s_k\sqrt{[n]}}{2\sqrt{1-c_{n,k}^2}}\label{edge}
\end{equation}
for $E,\B,\PP$, respectively; reverse edges are conjugates.
Finally the exact Hamiltonian boost is
\begin{equation}
 \Bt=\frac{(a_q+a_q^\dagger)\otimes(1-R)
                   -1\otimes(b_q+b_q^\dagger)}{2\sqrt\eps}.
 \label{boost}
\end{equation}

\newpage
\section{Proof of the finite-word limit and operator consequence}
Write $\zeta_j=D(-\sqrt m)\ket j$.
The displacement identity gives
\[
 \|(N/m-1)\zeta_j\|^2=(2j+1)/m+j^2/m^2.
\]
Its coefficients are the Poisson amplitude times a fixed-degree
Charlier polynomial. Explicitly, for the first three levels,
\begin{align*}
 (\zeta_0)_n&=(-1)^n e^{-m/2}m^{n/2}/\sqrt{n!},\\
 (\zeta_1)_n&=(m-n)(\zeta_0)_n/\sqrt m,\\
 (\zeta_2)_n&=((n-m)^2-n)(\zeta_0)_n/(\sqrt2m).
\end{align*}
For fixed $j$, the squared polynomial is bounded by
$C_jm^j(1+n/m)^{2j}$ for $m\ge1$.

Let $S_-\ket n=\ket{n-1}$ for $n\ge1$ and $S_+=S_-^\dagger$.
Comparison to $a_1/\sqrt m$, followed by displacement, gives
\[
 \|(S_-+1)\zeta_j\|\le
 \|(N/m-1)\zeta_j\|+\sqrt{j/m}\longrightarrow0.
\]
The analogous bound for $S_+$ includes the vanishing vacuum
boundary term. Every fixed string of shifts therefore acts as
its parity sign on $\zeta_j$, up to norm error tending to zero.

On $m/2\le n\le2m$ and bounded $k$, the diagonal coefficients
above converge uniformly to $d+k,0,0$.
The $E$ edge vanishes, while the $\B$ and $\PP$ edges tend to
$-\sqrt{(k+1)(2d+k)}/2$ and
$-i\sqrt{(k+1)(2d+k)}/2$. The minus sign from $S_-$
turns these into $B_0,P_0$. In \eqref{boost}, the common
oscillator edges are $O_{d,k}(\sqrt{\eps m})$, while
$s_k/\sqrt\eps\to\sqrt{(k+1)(2d+k)}$; hence $\Bt\to-B_0$.

For a word of fixed length $L$, only $|\Delta n|,|\Delta k|\le L$
occur. From $[n]\le n$ and $1-e^{-x}\ge xe^{-x}$, each
coefficient has a bound of the form
$C_{d,k,L}(1+\eps(n+L))^L e^{L\eps(n+k+d+L)}$ after
enlarging the constant. Squaring this bound and using the
polynomial-times-Poisson estimate controls the omitted tails:
exponential tilting changes the mean to $me^{O_L(\eps)}$,
and the prefactor is $e^{m(e^{O_L(\eps)}-1)}\to1$.
Outside $[m/2,2m]$, exponential concentration beats every
fixed polynomial. The finite path expansion of the word now
proves \eqref{wordlimit}. This argument controls products
directly, rather than inferring them from convergence of factors.

\paragraph{Strong resolvents.}
Each finite-$q$ LS combination is the self-adjoint direct sum
of its finite total-level blocks; $\Bt$ is bounded at fixed $q$.
The limiting Jacobi operators are essentially self-adjoint on
finite support since their off-diagonal coefficients grow
linearly and satisfy the Carleman criterion. The same applies
to $Q_0=E_0-P_0$; $E_0$ is diagonal.
Core convergence implies strong-resolvent convergence for
the conjugated individual operators, including $Q=E-\PP$.
Indeed, for $w=(A_0-z)v$, $v\in\D$, the resolvent difference
is bounded by $|\operatorname{Im}z|^{-1}\|(A_0-A_{q,m})v\|$;
such $w$ form a dense set. The uniform resolvent bound extends
the result to the whole space. The corresponding unitary groups
converge strongly at each fixed time.

\newpage
\section{A positive null-generator candidate}
In the moving frame, define $Q_0=E_0-P_0$ and
$T_0=-B_0$, the limiting physical boost. Then
\begin{equation}
 [T_0,Q_0]=iQ_0,\qquad Q_0\ge0.\label{null}
\end{equation}
The commutator follows from the displayed algebra. Positivity
has a direct form proof: for finite-support $v=(v_k)$,
\begin{equation}
 \langle v,Q_0v\rangle=
 \frac12\sum_{k\ge0}
 \left|\sqrt{2d+k}\,v_k+i\sqrt{k+1}\,v_{k+1}\right|^2.
 \label{positive}
\end{equation}
Its closure is nonnegative. A zero eigenvector would obey
$v_{k+1}=i\sqrt{(2d+k)/(k+1)}v_k$, whose squared coefficients
are proportional to $(2d)_k/k!$. They are not summable for
$d>0$, so zero is not an eigenvalue.

The phase transformation $v_k=i^kw_k$ turns $Q_0$ into the
Laguerre Jacobi matrix with diagonal $d+k$ and off-diagonal
$-\sqrt{(k+1)(2d+k)}/2$. The recurrence and orthogonality of
$L_k^{(2d-1)}(2x)$ \cite{DLMF} give the seed spectral measure
\begin{equation}
 d\nu_d(x)=\frac{2^{2d}}{\Gamma(2d)}x^{2d-1}e^{-2x}\,dx,
 \qquad x>0.\label{gamma}
\end{equation}
Thus $\langle Q_0\rangle=d$, $\operatorname{Var}(Q_0)=d/2$,
and $\langle e^{-aQ_0}\rangle=(2/(2+a))^{2d}$ for $a\ge0$.
This is the standard positive discrete-series realization already
used as a kinematic model in the project's earlier translation
note. The present step identifies its finite-word and resolvent
approximation by actual chord generators in a specified moving
frame; it does not create a new Laguerre representation.

In particular, Proposition~\ref{words} proves
\begin{equation}
 \|([\Bt,Q]-iQ)\U v\|\longrightarrow0,\qquad v\in\D.
\end{equation}
The finite-$q$ matrix $Q$ itself need not be positive.
Equation \eqref{null} belongs to its self-adjoint limiting
operator. The associated boost action rescales $Q_0$ as
$e^{i\eta T_0}Q_0e^{-i\eta T_0}=e^{-\eta}Q_0$.

\paragraph{The remaining translation problem.}
A single positive dilation-covariant operator is a null
translation candidate. It is not a complete Poincare dictionary.
The other LS combination $E_0+P_0$ does not commute with $Q_0$,
so the two cannot simply be called the two light-cone momenta.
The earlier construction $P^+=\mu Q_0$,
$P^-=M^2/(\mu Q_0)$ still inserts $M$ and $\mu$ by hand.
Determining the opposite null action and its multiparticle
composition from the microscopic model, rather than selecting
an arbitrary inverse, is the next physical question. In particular
$d=2$ remains a representation weight, not a mass determination.

\newpage
\section{An insertion of the full selected density}
The established covariant algebra has the self-adjoint fields
$h,U$ and the actual density
\begin{equation}
 \rho=a+s(h-W)^2,\quad
 s=\sqrt q,\quad W=\sqrt\lambda U,\quad
 a=1-s(1+\lambda)>0.\label{rho}
\end{equation}
Here $U$ is a unit q-Gaussian. Its operator norm gives the
uniform estimate
\begin{equation}
 \|W\|\le\sqrt{\frac{\eps}{1-e^{-\eps}}}=O(1).
 \label{Wnorm}
\end{equation}
This bound retains $U$ on the full tracial Hilbert space; no
energy-only replacement of $\rho$ is made.

Specify the preparation being compared. Embed the one-matter
module in the full tracial representation, choose a unit
$v\in\D$, and let $\zeta=\U v$. Apply the density on the left:
\begin{equation}
 \zeta_\rho=\frac{L_{\sqrt\rho}\zeta}
                    {\|L_{\sqrt\rho}\zeta\|}.
 \label{protocol}
\end{equation}
This is a vector-insertion protocol. It is not automatically
the original nonstationary functional
$\tau(\rho X^\dagger\alpha_t(X))$, whose density ordering
must be retained separately.

\begin{proposition}\label{selected}
In the window of Proposition~\ref{words}, set
\[
 \eta=\frac{\|(h+2\sqrt m)\zeta\|+\|W\|}{2\sqrt m}
               +\frac{\sqrt a}{2\sqrt{s m}}.
\]
Then
\begin{align}
 \left\|\left(\frac{\sqrt\rho}{2\sqrt{s m}}-1\right)
                  \zeta\right\|&\le\eta\longrightarrow0,\\
 \|\zeta_\rho-\zeta\|&\le2\eta.\label{selectedbound}
\end{align}
For every bounded operator $A$ in this common representation,
the expectation difference is at most $4\eta\|A\|$.
The squared insertion norm $\|L_{\sqrt\rho}\zeta\|^2$
divided by $4sm$ tends to one.
\end{proposition}
\begin{proof}
On the first factor $h=a_q+a_q^\dagger$.
The inequality $0\le1-\sqrt{[n]/n}\le\eps n/2$ bounds the
q-oscillator error on $\zeta_j$ by
$O_j(\eps m^{3/2})$. The ordinary displaced Hamiltonian
has residual norm $\sqrt{2j+1}$ about $-2\sqrt m$.
Hence $\|(h+2\sqrt m)\zeta\|/\sqrt m\to0$ for finite $v$.
Put $A_*=h-W$. Spectral calculus gives
$\|(|A_*|-2\sqrt m)\zeta\|\le\|(A_*+2\sqrt m)\zeta\|$
and
$\|\sqrt{a+sA_*^2}-\sqrt{s}|A_*|\|\le\sqrt a$.
Together with \eqref{Wnorm} these prove the first bound.
Normalizing a vector at distance at most $\eta$ from a unit
vector changes that distance to at most $2\eta$.
The bounded-observable estimate follows by two applications
of Cauchy--Schwarz.
\end{proof}

For the true endpoint heat vector, the preceding stage supplies
the same Hamiltonian concentration under the stronger window
$\eps m^2\to0$, so this bound also applies there.
It proves that this thermal insertion loses leading preparation
information. Recovering the original selected response requires
its ordering and subleading terms, not just this scalarized limit.

\newpage
\section{Verification and next calculation}
The exact checks use $q=1/4,4/9$, $d=1,2,3$ and total levels
zero through seven: 1,224 Gram equalities and 864 ladder
actions compare \eqref{frame} with the original nonorthogonal
LS representation. Twenty-one symbolic Laguerre recurrences
and three exact finite positivity forms check the spectral step.

For moving vectors we use $m=8,32,128,512$, $\eps=m^{-3}$,
$d=2$, $j=0,1,2$ and $k=0,1,2,3$ (48 cases).
Each case tests nine words, all three commutators, the Casimir,
and the boost match. The grid is padded beyond every applied
shift. Two common-oscillator cutoffs differ by 48.
The largest word refinement difference is $1.08\times10^{-23}$;
a Poisson polynomial envelope bounds the omitted input norm
squared by $3.18\times10^{-48}$.
These are distinct checks: the input tail alone does not certify
every unbounded word. Its analytic control is in the proof.

\begin{center}
\begin{tabular}{rrrrr}
\toprule
 & \multicolumn{2}{c}{$\|([\B,\PP]-iE)\U\ket{j,k}\|$}
 & \multicolumn{2}{c}{$\|([\Bt,Q]-iQ)\U\ket{j,k}\|$}\\
$m$ & $(0,0)$ & $(2,3)$ & $(0,0)$ & $(2,3)$\\
\midrule
8 & 0.479014 & 3.206959 & 0.337907 & 2.461123\\
32 & 0.123308 & 1.043909 & 0.114994 & 0.933263\\
128 & 0.031132 & 0.266054 & 0.047929 & 0.374348\\
512 & 0.007805 & 0.066461 & 0.022581 & 0.172697\\
\bottomrule
\end{tabular}
\end{center}
The errors on excited states are visibly larger than on the
seed; a seed check would not establish the theorem's scope.
Eight words on the largest-displacement $j=2,k=3$ vector were
recomputed by an independent sparse direct-coefficient routine
at 50 decimal digits. The largest relative vector difference is
$1.40\times10^{-13}$. The $[\B,\PP]$ cancellation ratio reaches
$733.45$, and the null-commutator ratio reaches $376.62$.
They are retained in the output. Floating-point diagnostics are
not interval certificates. The other two commutators have
cancellation ratios up to $4.03\times10^8$ and $1.52\times10^9$;
their tiny residuals are precision-sensitive and are not covered
by the eight-word high-precision comparison. Their convergence
rests on Proposition~\ref{words}. No successful earlier scientific
driver was rerun or its output modified.

The next calculation should use the established multiparticle
chord coproduct to determine whether the positive candidate
$Q_0$ has the required additive action and what fixes its
opposite null partner. In parallel, the insertion ordering in
the actual selected correlator must be carried through this
frame; its $U$ contribution cannot be recovered from the leading
scalarization alone. Four-bilocal and finite-$G$ dynamics remain
open. This note advances the microscopic generator map rather
than declaring those problems solved.

\begin{thebibliography}{9}
\bibitem{LS} H.~W. Lin and D. Stanford,
\emph{A symmetry algebra in double-scaled SYK},
\href{https://arxiv.org/abs/2307.15725}{arXiv:2307.15725},
Secs.~2--4; conventions from (18), (45)--(48), (147).
\bibitem{Oku} K. Okuyama, \emph{More on doubled Hilbert space
in double-scaled SYK},
\href{https://arxiv.org/abs/2404.02833}{arXiv:2404.02833}, Sec.~3.
\bibitem{Xu} J. Xu, \emph{On Chord Dynamics and Complexity Growth
in Double-Scaled SYK},
\href{https://arxiv.org/abs/2411.04251v3}{arXiv:2411.04251v3}, Sec.~3.2.
\bibitem{DLMF} NIST DLMF,
\href{https://dlmf.nist.gov/18.3}{18.3} and
\href{https://dlmf.nist.gov/18.9}{18.9}, Laguerre orthogonality and recurrence.
\bibitem{DLMFq} NIST DLMF,
\href{https://dlmf.nist.gov/18.28}{18.28(iii)}, Al-Salam--Chihara polynomials.
\end{thebibliography}
\end{document}
